\documentclass[10pt,a4paper]{amsart}
\usepackage[margin=19mm]{geometry}
\usepackage{amsmath,amssymb,mathtools}
\usepackage[scr=rsfs,cal=euler]{mathalfa}
\usepackage{xfrac}
\usepackage[unicode,hidelinks,bookmarks=false]{hyperref}
\newcommand{\ie}{i.e.\ }
\newcommand{\cf}{cf.\ }
\newcommand{\resp}{respectively\ }
\newcommand{\Subsection}{\subsection}
\providecommand{\nobreakdash}{\nobreak\hbox{-}}
\setlength{\emergencystretch}{2em}
\multlinegap0pt
\usepackage{bm}
\usepackage{bbm}
\usepackage{dsfont}
\usepackage{xspace}
\usepackage{cancel}
\usepackage{mathtools}
\usepackage{amsthm}
\makeatletter
\@ifundefined{equation}{\newtheorem{equation}{Equation}}{}
\@ifundefined{thm}{\newtheorem{thm}[equation]{Theorem}}{}
\makeatother
\newcommand{\G}{\Gamma}
\newcommand{\Se}{\mathcal{S}}
\newcommand{\al}{\alpha}
\newcommand{\calS}{{\mathcal S}}\newcommand{\cals}{{\bf s}}\newcommand{\calst}{{\bf st}}\newcommand{\cS}{{\mathcal S}}
\newcommand{\oo}{\mathfrak{o}}
\newcommand{\w}{\omega}
\title[Flat semigroups and weighted homogeneous surface singulariti]{Flat semigroups and weighted homogeneous surface singularities}
\author{Zsolt Baja, Tam\'as L\'aszl\'o}
\date{Source: arXiv:2302.14675v2 (CC BY 4.0)}
\begin{document}
\maketitle

\noindent\emph{Source and licence.} Flat semigroups and weighted homogeneous surface singularities. Version arXiv:2302.14675v2 (CC BY 4.0), distributed under the Creative Commons Attribution 4.0 International licence, as recorded in the arXiv licence metadata for this exact version and shown on the abstract page https://arxiv.org/abs/2302.14675v2. Full rights notice: licenses/arxiv-2302.14675-CC-BY-4.0.txt.\par\smallskip

\noindent\textit{Editorial scope.} Counted unit: the Thm whose source label is \texttt{flat:divisor} in the article (source0.tex lines 1357--1368, source label \texttt{flat:divisor}), delivered together with the article's own complete original proof of it (source0.tex lines 1369--1427, one \texttt{proof} environment of 59 lines). No mathematics is written, repaired or re-derived: statement and proof are the article's own text line for line. Required context retained uncounted: eq:w' (lines 445--447); thm\_flat\_rep (lines 1061--1064). Every definition, display and label of those lines is retained unchanged. Omitted from this package and not redistributed: 1--444, 448--1060, 1065--1356, 1428--1845. Nothing inside a retained excerpt is omitted or altered: every retained body is delivered exactly as the article's lines are written, so no omission receipt of any kind accompanies this package. Citations: the retained excerpts carry no citation command at all, so no bibliography entry of the article is transcribed and no reference is redistributed. Reference adaptation: none was needed and none was made; every reference inside the delivered unit resolves inside it, so no unresolved reference or citation is produced. Printed slips kept literal: no printed word, symbol, hypothesis or number of the counted statement or of its proof was altered. Layout and class: the article's own class is \texttt{amsart} with packages including amsfonts, amsthm, amsmath, amssymb, amscd, inputenc, geometry, enumerate, color, soul, array, tabularx, xypic, epic; the wrapper is the lane's pinned \texttt{amsart} editorial wrapper at 10pt on A4, which restates the theorem-like environments the retained text needs and adds the article's own macro definitions that the retained text needs (\texttt{\textbackslash G}, \texttt{\textbackslash Se}, \texttt{\textbackslash al}, \texttt{\textbackslash calS}, \texttt{\textbackslash oo}, \texttt{\textbackslash w}), with extra wrapper packages (bm, bbm, dsfont, xspace, cancel, mathtools). The delivered numbering is the unit's own. Assets: no retained span contains a figure environment, a graphics-inclusion command, a PDF-page inclusion, a table or a TikZ picture; no image file is copied into this package, the package declares no asset dependency and no image file is redistributed with it. Licence: version arXiv:2302.14675v2 of this article is distributed under the Creative Commons Attribution 4.0 International licence, as recorded in the arXiv licence metadata for this exact version and shown on the abstract page https://arxiv.org/abs/2302.14675v2. A full rights notice is delivered at licenses/arxiv-2302.14675-CC-BY-4.0.txt. Characters: every retained excerpt is pure ASCII, so the delivered engine, XeLaTeX with its default Latin Modern text font, reports no missing character.
\begin{equation}\label{eq:w'}
\omega_i\omega_i'\equiv 1 \, (\mathrm{mod} \,\alpha_i), \ \ 0< \omega_i'<\alpha_i.
\end{equation}

\begin{thm}
    \label{thm_flat_rep}
    Every flat semigroup is representable.
\end{thm}

\begin{thm}
     \label{flat:divisor}
     Let $\G$ be an SSR  graph defined by the Seifert invariants   
     \[
         Sf=\big(-b_0,s_1\times(\al_1,\w_1),\dots,s_n\times(\al_n,\w_n)
         \big) .    
     \]
     If the numbers $\al_i\geq 2$  are pairwise relatively prime integers and $\gcd(\al_i,s_i)=1$ for every  $i$,
     then its associated semigroup is a quotient of a flat semigroup.
     In fact,
     $\Se_{\G}=G(\al,s_1 \widehat{\al}_1,\dots,s_n\widehat{\al}_n)/\oo$.
 \end{thm}

 \begin{proof}
 First of all we observe that $\gcd(\al_i,\oo)\neq 1$ implies $\gcd(\al_i,s_i)\neq 1$, hence by the assumptions we must have $\gcd(\al_i,\oo)=1$ for every $i\in\{1,\dots,n\}$. Indeed, for a fixed $i$ the expression $\oo=\al b_0-\sum_{j\neq i}s_j\w_j\widehat{\al}_j - s_i\w_i\widehat{\al}_i$ implies the identity $s_i\w_i\widehat{\al}_i\equiv 0\,(\mathrm{mod }\,{\gcd(\al_i,\oo)})$, which simplifies to $s_i\w_i\equiv 0\,(\mathrm{mod }\,{\gcd(\al_i,\oo)})$, since $\gcd(\al_i,\al_j)=1$ for $j\neq i$. Then, by multiplying the last congruence with $\w_i'$ and using (\ref{eq:w'}) yields that $s_i\equiv 0\,(\mathrm{mod }\,{\gcd(\al_i,\oo)})$, which supports our claim. 
 
 
    Now, using similar ideas as in the proof of the Theorem \ref{thm_flat_rep} we consider the non-negative integers $k_1,\dots,k_n$ as the solutions of the equations
     \begin{equation}\label{eq:cong1}
      k_1 \al_1+\w_1\equiv
     \dots\equiv k_n \al_n+\w_n\equiv 0\,(\mathrm{mod }\,{\oo}),    
     \end{equation}
     with $0\leq k_i< \oo$ for every $1\le i\le n$. 
     These, on one hand, allow us to define the numbers 
     \[
        \tilde{\w}_i=\frac{1}{\oo} (k_i\al_i+\w_i)\quad
        \mbox{ for every }  1\le i\le n,
     \]
     which satisfy $\gcd(\tilde{\w}_i,\al_i)=1$ and $0<\tilde{\w}_i<\al_i$. 
     On the other hand, they imply that $b_0+\sum_{i=1}^ns_ik_i \equiv 0\,(\mathrm{mod }\,{\oo})$ too, hence one can define the new central  decoration as $\tilde{b}_0=(b_0+\sum_{i=1}^ns_ik_i)/\oo$. Indeed, from the equations (\ref{eq:cong1}), one gets $\sum_{i=1}^n\al s_i k_i +s_i\w_i\widehat{\al}_i\equiv 0\,(\mathrm{mod }\,{\oo})$. Furthermore,  the definition of $\oo$ reads as $\oo=\al b_0-\sum_i s_i\w_i\widehat{\al}_i$ which,  applied to the previous equation,  gives $\al(b_0+\sum_i s_i k_i) \equiv 0\,(\mathrm{mod }\,{\oo})$.  Since $\gcd(\al,\oo)=1$ we deduce that  $(b_0+\sum_i s_i k_i) \equiv 0\,(\mathrm{mod }\,{\oo})$.
     
     We consider the SSR graph $\tilde{\G}$ defined by the Seifert invariants
     \[
       Sf=\big(-\tilde{b}_0,s_1\times(\al_1,\tilde{\w}_1),\dots,s_n\times(\al_n,\tilde{\w}_n)
    \big).     
    \]
    In the followings, any of the numerical data associated with $\tilde{\G}$
    will be distinguished by the \textquoteleft tilde\textquoteright\, notation, 
    eg.  $\tilde{e}$ will stand for the orbifold Euler number of $\tilde{\G}$. The first observation is that $\tilde{e}=e/o<0$, hence $\tilde{\G}$ is negative definite. Furthermore, $\tilde{\al}=\al$ and one implies  that  $\tilde{\oo}=1$. Then, by the assumptions of the theorem we conclude that $\tilde{\G}$ is the canonical representative of the flat semigroup $\calS_{\tilde{\Gamma}}=G(\al,s_1 \widehat{\al}_1,\dots, s_n\widehat{\al}_n)$. 
    
    On the other hand, if we look at the quasi-linear function  $\tilde{N}(\ell)=\tilde{b}_0 \ell -\sum_{i=1}^n s_i\lceil\tilde{\w}_i \ell/\al_i\rceil$ associated with $\calS_{\tilde{\Gamma}}$, one can see that 
    $ \tilde{N}^{(\oo)}(\ell)=\tilde{N}(\oo\ell)=N(\ell)$. Hence $\Se_{\G}=\Se_{\tilde{\G}}/\oo$ which finishes the proof. 
  

    
%      \[
%      \tilde{\oo}=\al\left(\tilde{b}_0 -\sum_{i=1}^n s_i\frac{\tilde{w}_1}{\al_1}\right)
%      = \frac{\al}{\oo}\left(b_0+\sum_{i=1}^n s_i k_i
%      - \sum_{i=1}^n s_i\left(k_i+\frac{\w_i}{\al_i}\right)\right) =\frac{\al}{\oo}|e|=1.
%      \]
     
%      the function $\tilde{N}(l)=\tilde{b}_0 \ell -\sum_{i=1}^n s_i\lceil\tilde{\w}_i \ell/\al_i\rceil$. 
%      By  elementary  arguments and  computations 
%      we get that  
%      
%   
%     Our first observation is that $\tilde{\al}=\al$,
%     hence we can conclude the following identities
%      \[
%      \tilde{\oo}=\al\left(\tilde{b}_0 -\sum_{i=1}^n s_i\frac{\tilde{w}_1}{\al_1}\right)
%      = \frac{\al}{\oo}\left(b_0+\sum_{i=1}^n s_i k_i
%      - \sum_{i=1}^n s_i\left(k_i+\frac{\w_i}{\al_i}\right)\right) =\frac{\al}{\oo}|e|=1.
%      \]
%    Since $\oo=1$ and the numbers $\{\al_i\}_{i=1}^{n}$ are pairwise
%     relatively primes, we have that the  graph $\tilde{\G}$
%     is the canonical representant of the flat semigroup $\Se:=G(\al,s_1 \widehat{\al}_1,\dots, s_n\widehat{\al}_n)$,
%     hence $\Se_{\tilde{\G}}=\Se$.
% 
%     From the construction of $\tilde{N}$ we have that  $N=\tilde{N}^{(\oo)}$,
%      hence $\Se_{\G}=\oo^{-1}\Se_{\tilde{\G}}=\oo^{-1} \Se$,
%      which deduces our claim.
 \end{proof}

\end{document}
